% GENERATED FILE. Edit canonical lessons, then run npm run generate:books.
% canonical-source-hash: fnv1a64:ed44cb2bc28a97e3
% canonical-lessons: SA-C80-crow, SA-C80-parrot, SA-C80-peacock, SA-C80-tortoise, SA-C80-frog

\chapter{Crow, Parrot and Peacock}
\label{ch:sa-crow-parrot-peacock}
\hlchaptermodality{80}

\begin{chapteropening}
\textbf{By the end of this chapter:} \emph{I can name a crow, a parrot, a peacock, a tortoise and a frog.}

Complete the last lesson of chapter 80: i can name a crow, a parrot, a peacock, a tortoise and a frog.
\end{chapteropening}

\section[kākaḥ]{\sk{काकः} (kākaḥ) — a crow}

\label{lesson:SA-C80-crow}



\begin{quote}
Before the new one: say the Sanskrit for a mouse, then the Sanskrit for a cat.
\end{quote}

\subsection*{You'll want to know: \sk{काकः}}
\textbf{\sk{काकः}} (\emph{kākaḥ}) — \textquotedblleft{}a crow\textquotedblright{}, a black \textbf{\sk{विहगः}}. Say it and hear its call.

\begin{grammarlens}[title={What you've built}]
The bird that names itself.
\end{grammarlens}

\subsection*{Your turn}
\begin{itemize}
\raggedright
  \item \emph{Say it:} \emph{kākaḥ}
  \item \emph{Say it:} \emph{kākaḥ}, once more
  \item \emph{Say it:} say \emph{kākaḥ}
  \item \emph{Recall:} say the Sanskrit for a mouse, then the Sanskrit for a cat, then say \emph{kākaḥ} again
\end{itemize}

\begin{tcolorbox}[breakable,colback=teal!4,colframe=teal!35!black,title={Before you move on}]
What does \sk{काकः} mean? (\textquotedblleft{}A crow\textquotedblright{}.) Say it once more.
\end{tcolorbox}
\section[śukaḥ]{\sk{शुकः} (śukaḥ) — a parrot}

\label{lesson:SA-C80-parrot}



\begin{quote}
Before the new one: say the Sanskrit for a cat, then the Sanskrit for a crow.
\end{quote}

\subsection*{You'll want to know: \sk{शुकः}}
\textbf{\sk{शुकः}} (\emph{śukaḥ}) — \textquotedblleft{}a parrot\textquotedblright{}, the \textbf{\sk{हरितः}} bird that repeats words.

\begin{grammarlens}[title={What you've built}]
The green bird.
\end{grammarlens}

\subsection*{Your turn}
\begin{itemize}
\raggedright
  \item \emph{Say it:} \emph{śukaḥ}
  \item \emph{Say it:} \emph{śukaḥ}, once more
  \item \emph{Say it:} say \emph{śukaḥ}
  \item \emph{Recall:} say the Sanskrit for a cat, then the Sanskrit for a crow, then say \emph{śukaḥ} again
\end{itemize}

\begin{tcolorbox}[breakable,colback=teal!4,colframe=teal!35!black,title={Before you move on}]
What does \sk{शुकः} mean? (\textquotedblleft{}A parrot\textquotedblright{}.) Say it once more.
\end{tcolorbox}
\section[mayūraḥ]{\sk{मयूरः} (mayūraḥ) — a peacock}

\label{lesson:SA-C80-peacock}



\begin{quote}
Before the new one: say the Sanskrit for a crow, then the Sanskrit for a parrot.
\end{quote}

\subsection*{You'll want to know: \sk{मयूरः}}
\textbf{\sk{मयूरः}} (\emph{mayūraḥ}) — \textquotedblleft{}a peacock\textquotedblright{}, who dances when the \textbf{\sk{वर्षा}} comes.

\begin{grammarlens}[title={What you've built}]
The dancing bird.
\end{grammarlens}

\subsection*{Your turn}
\begin{itemize}
\raggedright
  \item \emph{Say it:} \emph{mayūraḥ}
  \item \emph{Say it:} \emph{mayūraḥ}, once more
  \item \emph{Say it:} say \emph{mayūraḥ}
  \item \emph{Recall:} say the Sanskrit for a crow, then the Sanskrit for a parrot, then say \emph{mayūraḥ} again
\end{itemize}

\begin{tcolorbox}[breakable,colback=teal!4,colframe=teal!35!black,title={Before you move on}]
What does \sk{मयूरः} mean? (\textquotedblleft{}A peacock\textquotedblright{}.) Say it once more.
\end{tcolorbox}
\section[kūrmaḥ]{\sk{कूर्मः} (kūrmaḥ) — a tortoise}

\label{lesson:SA-C80-tortoise}



\begin{quote}
Before the new one: say the Sanskrit for a parrot, then the Sanskrit for a peacock.
\end{quote}

\subsection*{You'll want to know: \sk{कूर्मः}}
\textbf{\sk{कूर्मः}} (\emph{kūrmaḥ}) — \textquotedblleft{}a tortoise\textquotedblright{}, slow, with its house on its \textbf{\sk{पृष्ठम्}}.

\begin{grammarlens}[title={What you've built}]
The slow one.
\end{grammarlens}

\subsection*{Your turn}
\begin{itemize}
\raggedright
  \item \emph{Say it:} \emph{kūrmaḥ}
  \item \emph{Say it:} \emph{kūrmaḥ}, once more
  \item \emph{Say it:} say \emph{kūrmaḥ}
  \item \emph{Recall:} say the Sanskrit for a parrot, then the Sanskrit for a peacock, then say \emph{kūrmaḥ} again
\end{itemize}

\begin{tcolorbox}[breakable,colback=teal!4,colframe=teal!35!black,title={Before you move on}]
What does \sk{कूर्मः} mean? (\textquotedblleft{}A tortoise\textquotedblright{}.) Say it once more.
\end{tcolorbox}
\section[maṇḍūkaḥ]{\sk{मण्डूकः} (maṇḍūkaḥ) — a frog}

\label{lesson:SA-C80-frog}



\begin{quote}
Before the new one: say the Sanskrit for a peacock, then the Sanskrit for a tortoise.
\end{quote}

\subsection*{You'll want to know: \sk{मण्डूकः}}
\textbf{\sk{मण्डूकः}} (\emph{maṇḍūkaḥ}) — \textquotedblleft{}a frog\textquotedblright{}, loudest in \textbf{\sk{वर्षा}}.

\begin{grammarlens}[title={What you've built}]
The rain singer.
\end{grammarlens}

\subsection*{Your turn}
\begin{itemize}
\raggedright
  \item \emph{Say it:} \emph{maṇḍūkaḥ}
  \item \emph{Say it:} \emph{maṇḍūkaḥ}, once more
  \item \emph{Say it:} say \emph{maṇḍūkaḥ}
  \item \emph{Recall:} say the Sanskrit for a peacock, then the Sanskrit for a tortoise, then say \emph{maṇḍūkaḥ} again
\end{itemize}

\begin{tcolorbox}[breakable,colback=teal!4,colframe=teal!35!black,title={Before you move on}]
What does \sk{मण्डूकः} mean? (\textquotedblleft{}A frog\textquotedblright{}.) Say it once more.
\end{tcolorbox}
